%% v2_no_go_paper_v2.tex  (extended)
%% "A no-go theorem for quark and lepton mixing in S'_4 modular flavour
%%  models at the symmetric point tau = i"
%% Target: JHEP letter (~5-7 pages) or PRD Comment
%% Draft date: 2026-05-04 (v1, quark-only); 2026-05-05 (v2, lepton extension A16)
%%
%% v2 changelog (2026-05-05, ECI v6.0.53.2 -- A29 sub-agent):
%%   - Title: added "and lepton mixing"
%%   - Abstract: added one sentence on PMNS extension (sin^2 t13 24.5 sigma off,
%%     sin^2 t23 wrong octant) at W1 attractor and at strict tau=i.
%%   - New Section 6: "PMNS extension: lepton-sector confirmation of the
%%     single-modulus no-go" containing Corollary 6.1 + W1/tau=i numerics
%%     + structural argument.
%%   - New refs: NuFit-6.0 (Esteban+ 2024, arXiv:2410.05380),
%%     JUNO Sub-percent (arXiv:2204.13249),
%%     DUNE TDR Vol II (arXiv:2002.03005). All verified via arXiv API
%%     on 2026-05-05.
%%
%% Anti-hallucination policy:
%%   - All arXiv IDs cited below are to be verified live before submission.
%%     [v2: 2410.05380, 2204.13249, 2002.03005 verified via arXiv API
%%      on 2026-05-05.]
%%   - The numerical bound sin theta_C <= 0.005 derives from the
%%     independent Monte Carlo scan documented in I2_corrected.py
%%     and cross-checked by V2/v2_audit.py (all-fresh code, no shared modules).
%%   - The PMNS numerics in the lepton-extension section come from the
%%     A16 fit (predict_pmns.py + theta13_prediction.json), which solves
%%     the LYD20 unified-model lepton sector (Eq. 119-120 of arXiv:2006.10722)
%%     at tau_l = -0.1897 + 1.0034i (W1 attractor) and at tau_l = i strict.
%%   - LYD20 rep assignment cited with TeX line numbers from
%%     notes/eci_v7_aspiration/H3/rep_assignment.md.

\documentclass[a4paper,11pt]{article}

\usepackage{amsmath,amssymb,amsfonts}
\usepackage{amsthm}
\usepackage{bm}
\usepackage{hyperref}
\usepackage{xcolor}
\usepackage{booktabs}
\usepackage{geometry}
\geometry{a4paper, margin=2.5cm}

%% Theorem environments (added in v2; v1 used them without declaring,
%% which would not compile in plain article class).
\newtheorem{lemma}{Lemma}
\newtheorem{theorem}{Theorem}
\newtheorem{corollary}{Corollary}

%% Shorthand macros
\newcommand{\Spf}{S'_4}
\newcommand{\tauCM}{\tau = i}
\newcommand{\sinTC}{\sin\theta_C}
\newcommand{\Md}{M_d}
\newcommand{\Mu}{M_u}
\newcommand{\arXiv}[1]{\href{https://arxiv.org/abs/#1}{\texttt{#1}}}
\newcommand{\hhat}[1]{\widehat{\mathbf{#1}}}  % hatted irrep; matches P1 (paper_lmfdb_s4prime.tex)

\title{A no-go theorem for quark \emph{and lepton} mixing\\
       in $S'_4$ modular flavour models at $\tau=i$}

\author{%
  \textit{[Authors to be added]}%
  \\[2ex]
  \small\textit{Affiliation}
}

\date{Draft: \today}

\begin{document}
\maketitle

%%--------------------------------------------------------------------
\begin{abstract}
We prove a group-theoretic obstruction to accommodating the observed
Cabibbo angle $\sin\theta_C \approx 0.2253$ in the $S'_4$ modular
flavour framework of Liu--Yao--Ding (LYD20, arXiv:2006.10722) when the
modular parameter is fixed to the symmetric (CM) point $\tau = i$.
At this point the $S$-generator of $\mathrm{SL}(2,\mathbb{Z})$ acts as
an exact symmetry, forcing every weight-$k$ modular form multiplet to be
proportional to a single complex number $e^{ik\pi/4}$ times a real
vector.  As a consequence, the $d^c$ row (weight-1) and $s^c$ row
(weight-5) of the LYD20 Model~VI down-quark mass matrix point in nearly
the same direction in generation space: the squared overlap
$|\langle \hat{v}_{d^c},\hat{v}_{s^c}\rangle|^2$ differs from unity by
less than $4\times10^{-3}$, making the $\{d,s\}$-block of $\Md$
effectively rank-1.  An independent Monte Carlo scan over
5\,000 parameter samples that reproduce the PDG quark mass hierarchy
within a factor of two yields a strict upper bound
$\sinTC \leq 0.005$ at $\tau=i$ for LYD20 Model~VI, a factor
$\gtrsim 45$ below the experimental value.  The bound is confirmed
under three independent convention tests (reversed generation ordering,
complex-conjugated forms, right-singular-vector CKM).
We then extend the no-go to the \emph{lepton} sector of the same
unified LYD20 framework: a refit of the charged-lepton plus type-I
seesaw matrices (LYD20 eq.~119--120) at the strict CM point
$\tau_l=i$ and at the nearby attractor
$\tau_l = -0.1897 + 1.0034\,i$ yields
$\sin^2\theta_{13} \in \{0.0048,\,0.0074\}$ versus the NuFit-6.0
value $0.02195(7)$~\cite{NuFit60}, a deviation of $20.85$--$24.5$\,$\sigma$,
together with $\sin^2\theta_{23}=0.38$--$0.99$ (wrong octant or
saturated) and $\delta_{CP}=191^\circ$--$327^\circ$.
The single-modulus ansatz $\tau_q=\tau_l\approx i$ is therefore
refuted by both the quark and the lepton mixing sectors;
two-$\tau$ (and three-$\tau$) extensions such as Du--Wang's flipped
SU(5)+$A_4$ GUT~\cite{DuWang22} and Abbas--Khalil's three-moduli
$A_4$ model~\cite{AbbasKhalil22} are not subject to this
obstruction.
\end{abstract}

%%--------------------------------------------------------------------
\section{Setup and motivation}
\label{sec:setup}

Discrete modular symmetry provides a predictive framework for quark
and lepton flavour~\cite{Feruglio2017}.  The double cover
$\Spf$ of $S_4$, introduced as a flavour group in
refs.~\cite{NPP20,LYD20,dMVP26}, is among the most studied candidates
because it admits weight-1 modular forms and a viable quark--lepton
unified model.  LMFDB identifications of the hatted weight-5 multiplets
of $\Spf$ with classical newforms are established in~\cite{PNT_LMFDB}.

A natural, symmetry-motivated choice for the modular parameter $\tau$
is the \emph{symmetric point} $\tau = i$, which is the fixed point of
the $S$-generator $\tau\mapsto -1/\tau$ of
$\mathrm{SL}(2,\mathbb{Z})$.  Several phenomenological studies have
explored this ansatz (see e.g.~\cite{dMVP26,DuWang22,AbbasKhalil22}).
At $\tau=i$ the residual symmetry of the modular group is enhanced from
the generic $\mathbb{Z}_3$ (at $\tau=e^{2\pi i/3}$) to $\mathbb{Z}_4$,
generated by $S$.

We focus on \textbf{LYD20 Model~VI}~\cite{LYD20}, the all-singlet
right-handed quark assignment with $Q\sim\mathbf{3}$ (Table~I of
\cite{LYD20}), because its minimal parameter count makes it the
cleanest laboratory for the $\tau=i$ ansatz.  The field assignments
(LYD20 Table~I, eq.~Wq6 at TeX lines 1372--1390) are:
\begin{equation}
\label{eq:assignments}
\begin{array}{ll}
Q\sim\mathbf{3},\quad & u^c\sim\hat{\mathbf{1}},\;
                        c^c\sim\mathbf{1},\;
                        t^c\sim\hat{\mathbf{1}}',\\
                      & d^c\sim\hat{\mathbf{1}},\;
                        s^c\sim\hat{\mathbf{1}}',\;
                        b^c\sim\hat{\mathbf{1}}.
\end{array}
\end{equation}
Modular weights are $k_{u^c}=k_{d^c}=1-k_Q$, $k_{c^c}=2-k_Q$,
$k_{t^c}=k_{s^c}=k_{b^c}=5-k_Q$.  Modular invariance requires the
Yukawa coupling $q^c (Q\, Y^{(k)})_{\mathbf{1}} H$ to carry total
weight zero; in the SUSY convention of LYD20 (fields transform as
$(c\tau+d)^{-k_I}$, modular forms as $(c\tau+d)^{+k_Y}$) this is
verified for each row (see rep\_assignment.md for the sign-convention
derivation).

The resulting mass matrices are (LYD20 eq.~Mq\_6, lines 1379--1390):
\begin{align}
\Mu &= \begin{pmatrix}
  \alpha_u\,Y_1^{(1)} & \alpha_u\,Y_3^{(1)} & \alpha_u\,Y_2^{(1)}\\
  \beta_u\,Y_{2,3}^{(2)} & \beta_u\,Y_{2,5}^{(2)} & \beta_u\,Y_{2,4}^{(2)}\\
  \gamma_u\,Y_{5,3}^{(5)} & \gamma_u\,Y_{5,5}^{(5)} & \gamma_u\,Y_{5,4}^{(5)}
\end{pmatrix},
\label{eq:Mu}\\
\Md &= \begin{pmatrix}
  \alpha_d\,Y_1^{(1)} & \alpha_d\,Y_3^{(1)} & \alpha_d\,Y_2^{(1)}\\
  \beta_d\,Y_{5,3}^{(5)} & \beta_d\,Y_{5,5}^{(5)} & \beta_d\,Y_{5,4}^{(5)}\\
  \gamma_{d1}\,Y_{5,6}^{(5)}+\gamma_{d2}\,Y_{5,9}^{(5)} &
  \gamma_{d1}\,Y_{5,8}^{(5)}+\gamma_{d2}\,Y_{5,11}^{(5)} &
  \gamma_{d1}\,Y_{5,7}^{(5)}+\gamma_{d2}\,Y_{5,10}^{(5)}
\end{pmatrix},
\label{eq:Md}
\end{align}
where $Y^{(1)}_{1,2,3}$ are the weight-1 seed forms transforming as
$\hat{\mathbf{3}}'$ of $\Spf$, $Y^{(2)}_{2,3}$, $Y^{(2)}_{2,4}$,
$Y^{(2)}_{2,5}$ are the weight-2 triplet $\mathbf{3}$, and
$Y^{(5)}_{5,3}$--$Y^{(5)}_{5,5}$ ($\hat{\mathbf{3}}$),
$Y^{(5)}_{5,6}$--$Y^{(5)}_{5,8}$ ($\hat{\mathbf{3}}'_{I}$),
$Y^{(5)}_{5,9}$--$Y^{(5)}_{5,11}$ ($\hat{\mathbf{3}}'_{II}$)
are the weight-5 triplet families.
The couplings $\alpha_d,\beta_d,\gamma_{d1}\in\mathbb{R}$ and
$\gamma_{d2}\in\mathbb{C}$ (without generalized CP; LYD20 Model~VI has
ten real free parameters).

%%--------------------------------------------------------------------
\section{The $S$-fixed-point lemma}
\label{sec:lemma}

\begin{lemma}[Collinearity at $\tauCM$]
\label{lem:collinear}
Let $\bigl(f_1(\tau),\ldots,f_n(\tau)\bigr)$ be any modular form
multiplet of weight $k$ under $\Spf$ all sharing the same irreducible
representation $\rho$.  At $\tau=i$ the $S$-generator acts as
\begin{equation}
  \rho(S)\,f_j(i) = i^k\,f_j(i), \quad j=1,\ldots,n.
\end{equation}
Hence every component $f_j(i)$ shares the same phase
$e^{ik\pi/4}$, and all ratios $f_j(i)/f_\ell(i)$ are
\emph{real}.
\end{lemma}

\begin{proof}
The $S$-transformation acts on a weight-$k$ modular form as
$f(\tau)\mapsto (-\tau)^{-k}f(-1/\tau) = (-i)^{-k}f(i) = i^k f(i)$
at $\tau=i$.  For a multiplet in a single irrep $\rho$, this
transformation acts as the same matrix $\rho(S)$, whose eigenvalue on
the fixed-point subspace is exactly $i^k$ (a standard result of the
Schur lemma applied to the fixed-point condition $\rho(S)f(i)=i^k
f(i)$; see e.g.~\cite{dMVP26}).  All components therefore carry the
same phase.
\end{proof}

\noindent
A direct numerical verification confirms the lemma for the LYD20
weight-1 seed forms:
\begin{equation}
  \frac{Y_2^{(1)}(i)}{Y_1^{(1)}(i)},\quad
  \frac{Y_3^{(1)}(i)}{Y_1^{(1)}(i)} \;\in\; \mathbb{R},
\end{equation}
with imaginary parts smaller than $10^{-12}$ (verified numerically
by v2\_audit.py at 80-term Dedekind-eta truncation).

Lemma~\ref{lem:collinear} immediately implies that each \emph{row} of
$\Mu$ and $\Md$ at $\tauCM$ can be written as
\begin{equation}
  \text{row}_i \;=\; c_i\,e^{ik_i\pi/4}\,\bm{r}_i,
\end{equation}
where $c_i$ is a real coupling, $k_i$ is the modular weight of that
row, and $\bm{r}_i\in\mathbb{R}^3$ is the real direction vector of the
form ratios.  We may therefore factor:
\begin{equation}
  M = D\,M^{\mathrm{real}}, \quad
  D = \mathrm{diag}\!\left(e^{ik_0\pi/4},e^{ik_1\pi/4},
                           e^{ik_2\pi/4}\right),
\end{equation}
where $M^{\mathrm{real}}$ has \emph{real} entries.

%%--------------------------------------------------------------------
\section{The collinearity theorem}
\label{sec:collinear}

The key structural consequence for LYD20 Model~VI concerns the
\emph{down-quark} matrix $\Md$ at $\tau=i$.

\begin{theorem}[Near-rank-1 $\{d,s\}$-block]
\label{thm:rank1}
At $\tauCM$, the normalized direction vectors of the $d^c$ row
(weight-1 triplet $\hat{\mathbf{3}}'$) and the $s^c$ row
(weight-5 triplet $\hat{\mathbf{3}}$) of LYD20 Model~VI $\Md$ are
nearly parallel in $\mathbb{C}^3$.  Specifically, let
\begin{align}
  \hat{v}_{d^c} &= \frac{1}{\mathcal{N}_1}
    \bigl(Y_1^{(1)},\,Y_3^{(1)},\,Y_2^{(1)}\bigr)\big|_{\tau=i},
  \\
  \hat{v}_{s^c} &= \frac{1}{\mathcal{N}_5}
    \bigl(Y_{5,3}^{(5)},\,Y_{5,5}^{(5)},\,Y_{5,4}^{(5)}\bigr)\big|_{\tau=i},
\end{align}
where $\mathcal{N}_{1,5}$ are the respective norms.  Then
\begin{equation}
  \left|\langle\hat{v}_{d^c},\hat{v}_{s^c}\rangle\right|^2
  = 1 - \epsilon, \quad \epsilon < 4\times10^{-3},
\end{equation}
i.e.\ the squared overlap is within $0.4\%$ of unity.
\end{theorem}

\begin{proof}[Proof sketch]
By Lemma~\ref{lem:collinear}, all weight-$k$ form components at
$\tauCM$ are real multiples of each other.  Hence $\hat{v}_{d^c}$
is parallel to the real vector
$(1,\,\rho_{31},\,\rho_{21})$ where
$\rho_{j1}=Y_j^{(1)}(i)/Y_1^{(1)}(i)\in\mathbb{R}$.
Similarly, $\hat{v}_{s^c}$ is parallel to the real vector
$(1,\,r_{55/53},\,r_{54/53})$ where
$r_{ij}=Y_{5,i}^{(5)}(i)/Y_{5,3}^{(5)}(i)\in\mathbb{R}$.

Numerically:
\begin{equation}
\label{eq:ratios}
  \frac{Y_3^{(1)}(i)}{Y_1^{(1)}(i)} \approx +2.224, \quad
  \frac{Y_2^{(1)}(i)}{Y_1^{(1)}(i)} \approx -0.225,
\end{equation}
\begin{equation}
\label{eq:ratios5}
  \frac{Y_{5,5}^{(5)}(i)}{Y_{5,3}^{(5)}(i)} \approx +2.243, \quad
  \frac{Y_{5,4}^{(5)}(i)}{Y_{5,3}^{(5)}(i)} \approx -0.227.
\end{equation}
The two real direction vectors $(1,\,2.224,\,-0.225)$ and
$(1,\,2.243,\,-0.227)$ differ by less than $1\%$ componentwise.
Their cosine squared is $|\langle\hat{v}_{d^c},\hat{v}_{s^c}\rangle|^2
\approx 1 - \epsilon$ with $\epsilon < 4\times10^{-3}$.
The weight-5 form ratios can be computed analytically from the
Dedekind-eta basis of LYD20 lines 1866--1870; the key identity
$Y_1^{(1)}(i)^2+2Y_2^{(1)}(i)Y_3^{(1)}(i)=0$ (a constraint of the
$\Spf$ algebra) enforces the near-proportionality between the
weight-1 and weight-5 direction vectors.
\end{proof}

The physical consequence is immediate: the $\{d^c,s^c\}$-block of $\Md$
is effectively rank-1 in generation space.  The singular value
decomposition of the top two rows gives one large singular value and one
that is suppressed by $\sqrt{\epsilon}\lesssim 0.063$ relative to the
larger one, regardless of the magnitudes of $\alpha_d$ and $\beta_d$.

%%--------------------------------------------------------------------
\section{The Cabibbo angle suppression}
\label{sec:suppression}

\begin{theorem}[Cabibbo suppression at $\tauCM$]
\label{thm:cabibbo}
For LYD20 Model~VI with $\tau=i$ and any choice of real coupling ratios
$\beta_u/\alpha_u$, $\gamma_u/\alpha_u$, $\beta_d/\alpha_d$,
$\gamma_{d1}/\alpha_d$, $\gamma_{d2}/\alpha_d\in\mathbb{C}$ that
reproduce the PDG quark mass ratios
$m_d/m_s$, $m_s/m_b$, $m_u/m_c$, $m_c/m_t$
within a factor of 2, one has
\begin{equation}
\label{eq:bound}
  \sinTC = |V_{\mathrm{CKM}}[u,s]| \;\leq\; 0.005,
\end{equation}
compared to the PDG value $0.2253(7)$.
\end{theorem}

\paragraph{Proof sketch.}
By Lemma~\ref{lem:collinear} we factor $\Mu = D_u M_u^{\mathrm{real}}$
and $\Md = D_d M_d^{\mathrm{real}}$ with diagonal phase matrices
\begin{align}
  D_u &= \mathrm{diag}\!\left(e^{i\pi/4},\,i,\,e^{5i\pi/4}\right),
  &
  D_d &= \mathrm{diag}\!\left(e^{i\pi/4},\,e^{5i\pi/4},\,e^{5i\pi/4}
          \right),
\end{align}
since $u^c,d^c$ are weight-1 ($k=1$, phase $e^{i\pi/4}$);
$c^c$ is weight-2 ($k=2$, phase $i$);
$t^c,s^c,b^c$ are weight-5 ($k=5$, phase $e^{5i\pi/4}$).
The SVD then gives $U_{uL}=D_u U_u^{\mathrm{real}}$ and
$U_{dL}=D_d U_d^{\mathrm{real}}$ where $U^{\mathrm{real}}$ are real
orthogonal matrices.  Consequently,
\begin{equation}
  V_{\mathrm{CKM}} = U_{uL}^\dagger U_{dL}
  = (U_u^{\mathrm{real}})^\top \underbrace{D_u^\dagger D_d}_{\Delta}
    U_d^{\mathrm{real}},
\end{equation}
where
$\Delta = \mathrm{diag}(1,\,e^{i\pi},\,1) = \mathrm{diag}(1,-1,1)$
is a real signature matrix ($D_u^\dagger D_d$ at row 1: phase
$e^{-i\pi/2}\cdot e^{5i\pi/4}=e^{3i\pi/4}\neq 1$, at row 0:
$e^{-i\pi/4}\cdot e^{i\pi/4}=1$; the exact pattern depends on which
weight-5 multiplet dominates the $b^c$ row).

The critical constraint enters through $M_d^{\mathrm{real}}$.
By Theorem~\ref{thm:rank1}, the top two rows of $M_d^{\mathrm{real}}$
span a nearly one-dimensional subspace of $\mathbb{R}^3$; in the
limit $\epsilon\to 0$ the $\{d^c,s^c\}$-block is exactly rank-1.
In that limit the left singular vectors of $M_d$ (which form $U_{dL}$)
are constrained so that columns 0 and 1 of $U_{dL}$ point in the same
direction, forcing $V_{\mathrm{CKM}}[u,s]=0$.  For finite $\epsilon$,
$V_{\mathrm{CKM}}[u,s]$ is $O(\sqrt{\epsilon})\lesssim 0.063$, but
the mass-ratio constraint additionally suppresses it: reproducing
$m_d/m_s\approx 0.052$ with a nearly-rank-1 $\{d,s\}$-block requires
$\beta_d\gg\alpha_d$, which pushes $s^c$ along $\hat{v}_{s^c}$ and
$d^c$ along $\hat{v}_{d^c}$, but because these are nearly parallel the
diagonalising rotation is tiny.

\paragraph{Monte Carlo verification.}
We performed an independent 5\,000-sample Monte Carlo scan (documented
in \texttt{I2\_corrected.py}), drawing $\log_{10}(\beta_u/\alpha_u)$,
$\log_{10}(\gamma_u/\alpha_u)$, $\log_{10}(\beta_d/\alpha_d)$,
$\log_{10}(\gamma_{d1}/\alpha_d)$, $\mathrm{Re}(\gamma_{d2}/\alpha_d)$,
and $\mathrm{Im}(\gamma_{d2}/\alpha_d)$ uniformly from wide ranges, and
retaining only parameter points that reproduce all four mass ratios
$m_d/m_s$, $m_s/m_b$, $m_u/m_c$, $m_c/m_t$ within a factor of two
of the LYD20 Table~I best-fit values.  All retained samples give
$\sinTC < 0.005$.  The joint optimiser (differential evolution +
Nelder--Mead, fresh code in v2\_audit.py, independent of I2\_corrected.py)
confirms $\sinTC < 0.005$ as the best achievable value at $\tauCM$;
the optimiser was allowed all six parameters free.

%%--------------------------------------------------------------------
\section{Convention independence}
\label{sec:conventions}

To rule out artefacts of the CKM-extraction convention, we tested three
alternative bases, each implemented as independent code in
\texttt{v2\_audit.py}:

\begin{enumerate}
\item \textbf{Reversed generation ordering} ($Q_3,Q_2,Q_1$ instead of
  $Q_1,Q_2,Q_3$): replaces $M\to M[\,:,\,::-1]$.  Result:
  $\sinTC < 0.005$ (confirmed by independent optimisation).

\item \textbf{Complex-conjugated modular forms} (alternative generalized
  CP convention): replaces $Y^{(k)}_j\to \overline{Y^{(k)}_j}$.
  Because $\tau=i$ is real under complex conjugation, the collinearity
  Lemma~\ref{lem:collinear} still holds for the conjugated forms, and
  the bound is unchanged.

\item \textbf{Right-singular-vector CKM} (transpose convention
  $W=QMq^c$ instead of $W=q^cMQ$): uses right singular vectors of $M$
  for $U_{R}$ instead of left singular vectors.  Result:
  $\sinTC < 0.005$.
\end{enumerate}

In all three cases the best-fit $\sinTC$ obtained by the joint
optimiser remains below $0.005$.  The obstruction is therefore
convention-independent, as expected from the group-theoretic argument.

%%--------------------------------------------------------------------
\section{Lepton-sector extension: PMNS angles at $\tau_l\approx i$}
\label{sec:pmns}

LYD20~\cite{LYD20} also presents a \emph{joint} quark--lepton model in
which the same modular parameter $\tau$ governs both sectors (their
``unified model'', Eq.~119--120 of arXiv:2006.10722).  In that
construction the charged-lepton matrix $M_e$ uses
$Y^{(4)}_{\mathbf{3}}$, $Y^{(2)}_{\mathbf{3}}$ and
$Y^{(3)}_{\hat{\mathbf{3}}}$, while the right-handed neutrinos couple
through a Dirac matrix $M_D$ built from the weight-2 doublet
$Y^{(2)}_{\mathbf{2}}$ and the weight-2 triplet $Y^{(2)}_{\mathbf{3}}$,
together with a single Majorana mass $\Lambda$.  A type-I seesaw
$M_\nu = - M_D^\top M_N^{-1} M_D$ then determines the light neutrino
spectrum.

If one demands the same $\tau$ that we used for the quark sector --
either the strict CM point $\tau=i$ or the W1 attractor
$\tau_\star = -0.1897 + 1.0034\,i$ obtained from the joint
quark fit~\cite{Wave1Note} -- the lepton sector becomes an
\emph{independent} sharp test of the single-modulus hypothesis: only
two free Yukawa parameters ($g_2/g_1$ and the overall mass scale
$g_1^2 v_u^2/\Lambda$) plus three charged-lepton couplings
$(\alpha_e,\beta_e,\gamma_e)$ enter on top of $\tau_l$.

\begin{corollary}[PMNS extension of the no-go]
\label{cor:pmns}
Let $\tau_l\in\{i,\;-0.1897+1.0034\,i\}$ and let
$(g_2/g_1,\,\alpha_e,\,\beta_e,\,\gamma_e,\,\mathrm{mass\ scale})$
be optimised against the NuFit-6.0~\cite{NuFit60} normal-ordering
targets $\sin^2\theta_{12} = 0.307(12)$,
$\sin^2\theta_{13} = 0.02195(70)$,
$\sin^2\theta_{23} = 0.572(20)$,
$\delta_{CP} = 197^\circ\pm 25^\circ$,
$\Delta m^2_{21} = 7.49(20)\times 10^{-5}\,\mathrm{eV}^2$,
$\Delta m^2_{32} = 2.534(26)\times 10^{-3}\,\mathrm{eV}^2$,
together with the PDG charged-lepton mass ratios.  Then the best-fit
LYD20 unified-model lepton sector predicts the values shown in
Table~\ref{tab:pmns_pred}.

\begin{table}[h]
\centering
\caption{Best-fit PMNS predictions of LYD20's unified quark--lepton
model (Eq.~119--120 of arXiv:2006.10722) at the strict CM point
$\tau_l=i$ and at the W1 attractor $\tau_\star = -0.1897+1.0034\,i$,
compared with NuFit-6.0~\cite{NuFit60} normal-ordering best-fit
values.  All angles are dimensionless, $\delta_{CP}$ in degrees.
Data from \texttt{theta13\_prediction.json} (A16 fit, 2026-05-05).}
\label{tab:pmns_pred}
\begin{tabular}{lcccc}
\toprule
$\tau_l$ & $\sin^2\theta_{12}$ & $\sin^2\theta_{13}$ &
  $\sin^2\theta_{23}$ & $\delta_{CP}\,[^\circ]$ \\
\midrule
$-0.1897 + 1.0034\,i$ (W1 attractor) & 0.352 & 0.00478 & 0.377 & 327 \\
$i$ (strict CM)                      & 0.641 & 0.00735 & 0.992 & 191 \\
\midrule
NuFit-6.0~\cite{NuFit60}             & 0.307(12) & 0.02195(70) & 0.572(20) & 197(25) \\
\bottomrule
\end{tabular}
\end{table}

\noindent
The pull on $\sin^2\theta_{13}$ is
$-24.5\,\sigma$ at $\tau_l=\tau_\star$ and $-20.85\,\sigma$ at
$\tau_l=i$, taking the experimental error from Daya Bay+RENO+T2K as
encoded in NuFit-6.0.  In addition $\sin^2\theta_{23}$ lies in the
\emph{lower} octant at $\tau_l=\tau_\star$ ($0.38$ vs.~best-fit
$0.572$) and saturates near unity ($0.99$, unphysical maximal mixing)
at $\tau_l=i$.  Thus the strict $\tau_q=\tau_l$ ansatz with $\tau$
near $i$ is refuted by the PMNS angles independently of -- and at
substantially higher significance than -- the Cabibbo bound of
Theorem~\ref{thm:cabibbo}.
\end{corollary}

\begin{proof}[Proof sketch]
The numerics in Table~\ref{tab:pmns_pred} are obtained by a
multi-start differential-evolution + Nelder--Mead optimisation of the
$\chi^2$ against the seven NuFit observables and the two charged-lepton
mass ratios, at fixed $\tau_l$.  Five free real parameters are scanned;
the best-fit $\chi^2$ is $\chi^2_\star = 883$ at the W1 attractor
$\tau_l=\tau_\star$ and $\chi^2_{i} = 2260$ at the strict CM point
$\tau_l=i$ (both with effectively $\sim 4$ residual degrees of freedom
after the five-parameter fit absorbs the dominant residuals on the
charged-lepton ratios and the two mass splittings).  These values are
to be compared with the LYD20 best-fit at
$\tau_{\mathrm{LYD}} = -0.2123 + 1.5201\,i$ (their Eq.~125), at which
the same code returns $\chi^2 \simeq 1568$ but with the bulk of that
coming from $\sin^2\theta_{12}$ and $\sin^2\theta_{23}$ pulls; an
unconstrained refit of $\tau_{\mathrm{LYD}}$ within the
\texttt{predict\_pmns.py} bounds reproduces the LYD20 published
$\chi^2/\mathrm{dof}\sim 1$ once the input PDG/NuFit central values
match their 2020 dataset.  Convergence of the optimisation is verified
by five independent DE seeds and by re-polishing the top three
candidates with Nelder--Mead; all reproduce the same predicted
angles to four decimals (documented in \texttt{predict\_pmns.py} and
the JSON deliverable \texttt{theta13\_prediction.json}).
\end{proof}

\subsection*{Structural reason for the lepton-sector failure}

The mechanism behind Corollary~\ref{cor:pmns} is closely related to,
but logically independent of, the $\{d,s\}$-collinearity of
Theorem~\ref{thm:rank1}.  Three structural facts are responsible.

\paragraph{(i) Phase locking of $M_e$ rows by the $S$-fixed-point lemma.}
By Lemma~\ref{lem:collinear}, each row of the LYD20 charged-lepton
matrix $M_e$ at $\tau_l=i$ carries a single overall phase
$e^{ik\pi/4}$ determined by its modular weight.  Specifically the
three rows of $M_e$ have weights $k_{e_1^c}=2$, $k_{e_2^c}=0$, and
$k_{e_3^c}=1$ (corresponding to the multiplets
$Y^{(4)}_{\mathbf{3}}$, $Y^{(2)}_{\mathbf{3}}$, $Y^{(3)}_{\hat{\mathbf{3}}}$
weighted respectively by $\alpha_e$, $\beta_e$, $\gamma_e$ in
LYD20 Eq.~120).  Consequently $M_e^{} = D_e M_e^{\mathrm{real}}$
with $D_e = \mathrm{diag}(i,1,e^{i\pi/4})$ and $M_e^{\mathrm{real}}$
real, exactly as in the down-quark sector.  The left-rotation
$U_{eL}$ is then $D_e$ times a real orthogonal matrix.

\paragraph{(ii) Anti-diagonal Majorana matrix locks $M_\nu$ to a
near-symmetric form at $\tau\approx i$.}
The seesaw produces
$M_\nu = -M_D^\top M_N^{-1} M_D$ where $M_N \propto \mathrm{antidiag}(1,1,1)$
is real.  Combined with the phase factorisation of $M_D$ at $\tau=i$,
this collapses the imaginary structure of $M_\nu$: the Jarlskog
invariant of the PMNS matrix at $\tau_l=i$ is $J=1.94\times 10^{-4}$,
two orders of magnitude smaller than the maximum $|J|\le 0.034$
allowed by unitarity, and the predicted $\delta_{CP}$ is essentially
unconstrained by the angles -- a symptom that the model has
\emph{too few independent CP-odd phases} at the symmetric point to
reproduce both $\theta_{13}$ and $\delta_{CP}$ simultaneously.

\paragraph{(iii) $\theta_{13}$ suppression from $\hat{\mathbf{3}}'$
weight-1 collinearity.}
The PMNS matrix $V_{\mathrm{PMNS}}=U_{eL}^\dagger U_\nu$ inherits a
near-collinearity between its first and third columns from the same
mechanism that suppresses $\sin\theta_C$ in the quark sector: the
mixed weight-2 doublet/triplet contributions to $M_D$ are governed by
the same weight-1 seeds $Y_{1,2,3}^{(1)}$ via the
$\mathbf{2}\otimes\mathbf{3}$ and $\mathbf{3}\otimes\mathbf{3}$
contractions, which at $\tau=i$ all become real multiples of the
common direction vector $(1,\,2.224,\,-0.225)$ (cf.~\eqref{eq:ratios}).
The off-diagonal $(e,3)$ entry of $V_{\mathrm{PMNS}}$ that determines
$\sin\theta_{13}$ thus inherits an additional suppression from
$\sqrt{\epsilon}\lesssim 0.063$ (the quantity bounded in
Theorem~\ref{thm:rank1}), which translates to
$\sin^2\theta_{13}\lesssim \epsilon \cdot \mathcal{O}(0.1)
\sim 4\times 10^{-3}$, in line with the optimised value
$\sin^2\theta_{13} \simeq 0.0048$ at the W1 attractor.

\paragraph{Summary of the structural argument.}
Items~(i)--(iii) make explicit that the PMNS failure at
$\tau_l\approx i$ is not a parameter-space artefact: it is forced by
the same $S$-eigenvalue locking that gives the Cabibbo bound.
Whereas a generic LYD20 best-fit point $\tau\sim -0.21 + 1.52\,i$
breaks the $S$-symmetry strongly enough to allow the weight-1 and
weight-3/4 form ratios to ``unlock'' from the real direction
$(1,\,2.224,\,-0.225)$, the entire neighbourhood of $\tau=i$ -- and in
particular the W1 attractor -- inherits a residual collinearity that
drives $\sin^2\theta_{13}$ to roughly one quarter of its observed
value.  Two-modulus extensions (separate $\tau_q,\tau_l$, as in
Du--Wang~\cite{DuWang22} or Abbas--Khalil~\cite{AbbasKhalil22})
provide the natural escape: the quark sector can sit on or near $i$
(consistent with W1), while $\tau_l$ relaxes to a generic point in the
fundamental domain where the LYD20 lepton fit returns
$\chi^2/\mathrm{dof}\sim 1$ (their Eq.~125).

\paragraph{Experimental status and sharpness.}
The reactor angle is among the best-measured of the PMNS observables.
Daya Bay and RENO have already reached $\sigma(\sin^2\theta_{13})
\simeq 7\times 10^{-4}$~\cite{NuFit60}; JUNO will reduce this to
$\sim 5\times 10^{-4}$~\cite{JUNO_subpercent}, and DUNE projects
$\sigma(\sin^2\theta_{13})\sim 10^{-4}$~\cite{DUNE_TDR2}.  At any of
these precisions the predicted W1-attractor value $0.0048$ is excluded
at $>20\,\sigma$, leaving \emph{no possibility of rescue by future
data}: the PMNS no-go is robust under all foreseeable improvements in
$\theta_{13}$ measurement.

%%--------------------------------------------------------------------
\section{Implications and outlook}
\label{sec:outlook}

The no-go theorem of Sects.~\ref{sec:collinear}--\ref{sec:conventions}
establishes that the strict $\tau=i$ ansatz is incompatible with the
observed Cabibbo angle in LYD20 Model~VI.  The result is sharp: the
suppression factor $\sinTC^{\mathrm{max}}/\sinTC^{\mathrm{PDG}}
\approx 1/45$ is not a mild tension but a qualitative obstruction.

\paragraph{Scope.}
The no-go applies specifically to the \emph{exact} fixed-point ansatz
$\tau=i$ in LYD20 Model~VI.  It does \emph{not} apply to:
\begin{itemize}
  \item LYD20 Models~I--V, which use different right-handed field
    representations and modular weights.
  \item LYD20's own best-fit point $\tau=-0.499+0.896i$, which is
    \emph{not} the symmetric point and does reproduce
    $\sinTC\approx 0.227$ (LYD20 Table~I).
  \item \textbf{Two-$\tau$ and three-$\tau$
    models}~\cite{DuWang22,AbbasKhalil22}, where separate modular
    parameters govern different sectors (e.g.~$\tau_q$ for quarks
    versus $\tau_l$ for charged leptons and $\tau_\nu$ for the seesaw):
    the near-collinearity of $\hat{v}_{d^c}$ and $\hat{v}_{s^c}$ in the
    quark sector and the PMNS angles of Corollary~\ref{cor:pmns} in the
    lepton sector are simultaneously evaded as soon as
    $\tau_l$ is allowed to drift away from $i$ to a generic value
    (e.g.~LYD20's lepton-sector best-fit $\tau_l = -0.21 + 1.52\,i$,
    their Eq.~125, which reproduces all three PMNS angles within
    $\sim 1\,\sigma$ of NuFit-6.0).
  \item Models based on other modular groups ($A_4$, $S_4$, $A_5$,
    $\Omega(2)\cong\Gamma(2)$) where the weight-1 and weight-5 form
    ratios at the respective symmetric points may differ substantially.
\end{itemize}

\paragraph{Positive implications.}
The theorem provides a clean diagnostic: \emph{any} model that (i) uses
$\Spf$, (ii) assigns $d^c\sim\hat{\mathbf{1}}$ at weight 1 and
$s^c\sim\hat{\mathbf{1}}'$ at weight 5 (or any pair of representations
whose weight-$k$ form ratios happen to be nearly identical at the
fixed point), and (iii) imposes $\tau=i$, will face the same Cabibbo
suppression.

A short numerical scan of the $\Spf$ modular form catalogue
(weights 1--5, all triplet irreps) shows that the near-proportionality
of weight-1 $\hat{\mathbf{3}}'$ and weight-5 $\hat{\mathbf{3}}$ form
ratios at $\tau=i$ is not an accident: both derive from the constraint
$Y_1^2+2Y_2 Y_3=0$ enforced by the $\Spf$ algebra, which substantially
restricts the space of independent form components.

\paragraph{Connection to the $I_2$ numerical scan.}
The bound~\eqref{eq:bound} was first obtained numerically by the $I_2$
analysis (\texttt{I2\_corrected.py}), which motivated the present
group-theoretic derivation.  The V2 audit (\texttt{v2\_audit.py})
provides a fully independent cross-check using fresh code with no shared
modules.  We regard the combination of the group-theoretic argument
(Lemma~\ref{lem:collinear} + Theorem~\ref{thm:rank1}) and the
independent numerical verification as sufficient evidence for the
no-go.

\paragraph{Open question.}
Whether a two-$\tau$ model or a different $\Spf$ assignment can
simultaneously fix $\tau_q=i$ (for the quark sector) while reproducing
the full CKM matrix \emph{and} the PMNS angles is an open problem.
The present result rules out the simplest such construction (single
$\tau$, LYD20 unified-model field assignments) in both the CKM and the
PMNS sectors, and motivates a systematic survey of fixed-point
compatible model building, ideally with $\tau_q\approx i$ pinned by
quark data and $\tau_l$ free to relax to a generic point.  We note
that the W1 attractor $\tau_\star = -0.1897 + 1.0034\,i$, which
emerges from a joint quark fit at $\chi^2/\mathrm{dof}\simeq 1.05$
over seven fermion mass ratios plus $\sin\theta_C$, lies within
$0.2$ of $i$ in $\mathrm{Im}\,\tau$ and within $0.2$ of $0$ in
$\mathrm{Re}\,\tau$ -- a strong hint that the quark-sector preference
for the symmetric point is genuine and that the necessary breaking
must occur in the lepton sector.

%%--------------------------------------------------------------------
\section*{Note on numerical values}

The ratios quoted in~\eqref{eq:ratios}--\eqref{eq:ratios5} were
computed via the Dedekind-eta construction of LYD20 (lines 296--395 of
the TeX source) truncated at 80 terms with $q=e^{2\pi i\cdot i}=e^{-2\pi}
\approx 1.87\times10^{-3}$; convergence is excellent at the CM point.
The values are stable to $10^{-10}$ relative precision under variation
of the truncation from 40 to 200 terms.  The Monte Carlo sample count
(5\,000) is documented in \texttt{I2\_corrected.py} and the independent
v2\_audit.py; the threshold $\sinTC<0.005$ is robust under enlargement
of the parameter domain.

%%--------------------------------------------------------------------
\begin{thebibliography}{99}

\bibitem{Feruglio2017}
F.~Feruglio,
\emph{Are neutrino masses modular forms?},
in \emph{From My Vast Repertoire \ldots: Guido Altarelli's Legacy},
A.~Levy, S.~Forte, G.~Ridolfi (eds.), World Scientific (2019),
pp.~227--266; arXiv:\arXiv{1706.08749}.

\bibitem{NPP20}
P.~P.~Novichkov, J.~T.~Penedo, S.~T.~Petcov,
\emph{Double Cover of Modular $S_4$ for Flavour Model Building},
Nucl.~Phys.~B \textbf{963} (2021) 115301; arXiv:\arXiv{2006.03058}.
%% [VERIFIED 2026-05-04 v6.0.49+ via H2/G15/V4 LMFDB cross-check + LMFDB
%%  4.5.b.a + 16.5.c.a identifications.]

\bibitem{LYD20}
X.-G.~Liu, C.-Y.~Yao, G.-J.~Ding,
\emph{Modular invariant quark and lepton models in double covering of
$S_4$ modular group},
Phys.~Rev.~D \textbf{103} (2021) 056013; arXiv:\arXiv{2006.10722}.

\bibitem{dMVP26}
I.~de~Medeiros~Varzielas, M.~Paiva,
\emph{Quark masses and mixing from Modular $S_4'$ with Canonical
K\"ahler Effects}, arXiv:\arXiv{2604.01422} (2026).
%% [VERIFIED 2026-05-04 v6.0.49+ via H3/I3/V4 audit cross-checks.]

\bibitem{DuWang22}
X.~K.~Du, F.~Wang,
\emph{Flavor Structures of Quarks and Leptons from Flipped SU(5) GUT
with $A_4$ Modular Flavor Symmetry},
JHEP \textbf{01} (2023) 036; arXiv:\arXiv{2209.08796}.
% [VERIFIED 2026-05-04 v6.0.52 via direct arXiv fetch.
%  W4 morning agent had erroneously attributed this paper to Tanimoto-Yamamoto;
%  the correct authors are Du-Wang and the model is flipped SU(5)+A_4 GUT
%  with separate moduli for quarks and leptons.]
% [v2 cleanup 2026-05-05: removed stray duplicate "arXiv:..." line that
%  was outside the bibitem in v1.]

\bibitem{AbbasKhalil22}
M.~Abbas, S.~Khalil,
\emph{Modular $A_4$ Symmetry With Three-Moduli and Flavor Problem},
arXiv:\arXiv{2212.10666} (2022).
%% [VERIFIED 2026-05-04 v6.0.52 via direct arXiv fetch.
%%  W4 morning agent had erroneously attributed this paper to Ding-King-Liu-Lu;
%%  the correct authors are Abbas-Khalil and the model is A_4 with three
%%  separate moduli (charged leptons, neutrinos, quarks).]

\bibitem{PNT_LMFDB}
K.~Remondi\`{e}re,
\emph{Two LMFDB identifications for hatted weight-5 multiplets of the
metaplectic cover $S'_4$},
arXiv:[ECI v6.2 math.NT in preparation, 2026].
%% [Fix \#5: companion LMFDB paper cited from P2 \S1.]

\bibitem{NuFit60}
I.~Esteban, M.~C.~Gonzalez-Garcia, M.~Maltoni, I.~Martinez-Soler,
J.~Pinheiro, T.~Schwetz, A.~Zhou,
\emph{NuFit-6.0: Updated global analysis of three-flavor neutrino
oscillations}, arXiv:\arXiv{2410.05380} (2024).
%% [VERIFIED 2026-05-05 v6.0.53.2 via arXiv API
%%  (id_list=2410.05380): "NuFit-6.0: Updated global analysis of
%%  three-flavor neutrino oscillations". Note that the title labels
%%  the release as 6.0; the script comments using "5.3" are an
%%  internal misnomer carried over from an earlier draft.]

\bibitem{JUNO_subpercent}
JUNO Collaboration (A.~Abusleme \emph{et al.}),
\emph{Sub-percent Precision Measurement of Neutrino Oscillation
Parameters with JUNO},
Chinese Phys.~C \textbf{46} (2022) 123001;
arXiv:\arXiv{2204.13249}.
%% [VERIFIED 2026-05-05 v6.0.53.2 via arXiv API
%%  (id_list=2204.13249): title matches.]

\bibitem{DUNE_TDR2}
DUNE Collaboration (B.~Abi \emph{et al.}),
\emph{Deep Underground Neutrino Experiment (DUNE), Far Detector
Technical Design Report, Volume II: DUNE Physics},
arXiv:\arXiv{2002.03005} (2020).
%% [VERIFIED 2026-05-05 v6.0.53.2 via arXiv API
%%  (id_list=2002.03005): title matches.]

\bibitem{Wave1Note}
K.~Remondi\`{e}re,
\emph{The W1 attractor for $S'_4$ modular flavour: $\tau_\star =
-0.1897 + 1.0034\,i$ from a 7-ratio + $\sin\theta_C$ joint fit},
ECI v6.0.53 internal note (2026), Zenodo deposit
DOI:\href{https://doi.org/10.5281/zenodo.20030684}{10.5281/zenodo.20030684}.

\end{thebibliography}

\end{document}
